\chapter{روش‌های کاهش بعد \texorpdfstring{\enparen{Dimensionality Reduction}}{(Dimensionality Reduction)}}
\label{ch:dimred}

\section{مقدمه و انگیزه کاهش بعد}
در عصر کلان‌داده، اشیاء غالباً با بردارهای بسیار طولانی در فضاهای با ابعاد نجومی (مثلاً $10^4$ تا $10^6$ بعد) مدل‌سازی می‌شوند؛ مانند ماتریس اسناد در برابر واژگان در پردازش متن، یا ماتریس کاربران در برابر کالاها در وب‌سایت‌های تجارت الکترونیک.

کاهش بعد\footnote{\lr{Dimensionality Reduction}} فرآیند نگاشت نقاط داده از فضای اصلی $d$-بعدی به فضایی با بعد بسیار کوچک‌تر $k \ll d$ است، به گونه‌ای که ساختار بنیادین، تنوع هندسی و فواصل میان نقاط داده تا حد ممکن حفظ گردد.

دلایل کلیدی کاهش بعد:
\begin{enumerate}
    \item \textbf{کاهش هزینه‌های محاسباتی و حافظه}: جایگزینی محاسبات سنگین در ابعاد بالا با عملیات سریع در ابعاد فشرده.
    \item \textbf{حذف نویز و افزونگی}: ابعاد داده معمولاً به شدت با یکدیگر همبسته هستند؛ کاهش بعد بخش‌های نویزی و بی‌اهمیت داده را فیلتر می‌کند.
    \item \textbf{استخراج معنا و مفاهیم پنهان}\footnote{\lr{Latent Concepts}}: ابعاد ترکیبی جدید اغلب مفاهیم کیفی معناداری را آشکار می‌سازند (مانند کشف سبک‌های ادبی یا ژانرهای سینمایی پنهان).
    \item \textbf{امکان مصورسازی}: کاهش ابعاد به ۲ یا ۳ بعد جهت رسم نمودارها و درک شهودی توسط انسان.
\end{enumerate}

\section{تجزیه مقادیر منفرد \texorpdfstring{\enparen{Singular Value Decomposition - SVD}}{(Singular Value Decomposition - SVD)}}
تجزیه مقادیر منفرد \enparen{SVD} یکی از شاهکارهای جبر خطی کاربردی است که بر روی هر ماتریس مستطیلی دلخواه قابل اعمال است.

\begin{theorem}[قضیه تجزیه مقادیر منفرد]
هر ماتریس حقیقی $A$ با ابعاد $m \times n$ و رتبه $r$ را می‌توان دقیقاً به حاصل‌ضرب سه ماتریس تجزیه کرد:
\begin{equation}
A = U \cdot \Sigma \cdot V^T
\end{equation}
که در آن:
\begin{itemize}
    \item $U$: ماتریس متعامد با ابعاد $m \times r$ که در شرط $U^T U = I$ صدق می‌کند. ستون‌های $U$، \textbf{بردارهای منفرد چپ}\footnote{\lr{Left Singular Vectors}} (بردارهای ویژه ماتریس متقارن $A A^T$) هستند.
    \item $\Sigma$: ماتریس قطری با ابعاد $r \times r$ با درایه‌های مثبت بر روی قطر اصلی:
    \begin{equation}
    \sigma_1 \ge \sigma_2 \ge \dots \ge \sigma_r > 0
    \end{equation}
    این مقادیر، \textbf{مقادیر منفرد}\footnote{\lr{Singular Values}} نامیده می‌شوند و برابر با جذر مقادیر ویژه ماتریس‌های $A^T A$ یا $A A^T$ هستند.
    \item $V$: ماتریس متعامد با ابعاد $n \times r$ که در شرط $V^T V = I$ صدق می‌کند. ستون‌های $V$، \textbf{بردارهای منفرد راست}\footnote{\lr{Right Singular Vectors}} (بردارهای ویژه ماتریس متقارن $A^T A$) هستند.
\end{itemize}
\end{theorem}

\subsection{تفسیر هندسی و فیزیکی SVD}
در تحلیل داده‌ها، ماتریس $A$ (مثلاً ماتریس کاربران در برابر فیلم‌ها) را می‌توان به صورت نگاشتی از فضای اشیاء به یک «فضای مفاهیم پنهان» تفسیر کرد:
\begin{itemize}
    \item هر بعد در فضای پنهان متناظر با یک سطر/ستون در $\Sigma$ است (مثلاً مفهوم ژانر علمی-تخیلی یا رمانس).
    \item سطر $i$-ام ماتریس $U$ بیانگر وابستگی کاربر $i$ به مفاهیم پنهان است.
    \item مقدار منفرد $\sigma_k$ نشان‌دهنده «قدرت و میزان واریانس پوشش داده‌شده توسط مفهوم $k$-ام» است.
    \item سطر $j$-ام ماتریس $V$ (ستون $j$-ام $V^T$) بیانگر میزان تعلق فیلم $j$ به مفاهیم پنهان است.
\end{itemize}

\subsection{تقریب با رتبه پایین و قضیه اکارت-یانگ-میرسکی}
اگر تنها $k$ مقدار منفرد نخست (بزرگ‌ترین مقادیر) را نگه داریم و سایر را صفر کنیم $(k < r)$:
\begin{equation}
A_k = U_k \cdot \Sigma_k \cdot V_k^T = \sum_{i=1}^k \sigma_i u_i v_i^T
\end{equation}

\begin{theorem}[قضیه اکارت-یانگ-میرسکی \enparen{Eckart-Young-Mirsky}]
ماتریس $A_k$، \textbf{بهترین تقریب ممکن با رتبه $k$} برای ماتریس $A$ در هر دو نرم فروبنیوس و نرم طیفی است:
\begin{equation}
\|A - A_k\|_F = \min_{\text{rank}(B) \le k} \|A - B\|_F = \sqrt{\sum_{i=k+1}^r \sigma_i^2}
\end{equation}
\end{theorem}
این قضیه تضمین می‌کند که خطای کاهش بعد توسط SVD حداقل ممکن در میان تمامی روش‌های خطی است.

\textbf{معیار حفظ انرژی برای انتخاب رتبه بهینه $k$}:
مجموع مجذور مقادیر منفرد ماتریس، بیانگر کل انرژی یا واریانس فروبنیوس موجود در داده‌ها است:
\begin{equation}
\text{Total Energy} = \|A\|_F^2 = \sum_{i=1}^r \sigma_i^2
\end{equation}
در مسائل داده‌کاوی واقعی، رتبه تقریب $k$ معمولاً به گونه‌ای برگزیده می‌شود که حداقل ۸۰ تا ۹۰ درصد از کل انرژی داده‌ها حفظ گردد:
\begin{equation}
\frac{\sum_{i=1}^k \sigma_i^2}{\sum_{i=1}^r \sigma_i^2} \ge 0.90
\end{equation}

\section{تحلیل مؤلفه‌های اصلی \texorpdfstring{\enparen{Principal Component Analysis - PCA}}{(Principal Component Analysis - PCA)}}
تحلیل مؤلفه‌های اصلی \enparen{PCA} به دنبال یافتن راستاهایی در فضا است که بیشترین پراکندگی و واریانس داده‌ها را در امتداد خود ثبت می‌کنند.

\begin{algorithmbox}[مراحل اجرای PCA]
\begin{enumerate}
    \item \textbf{مرکززدایی داده‌ها \enparen{Zero-Centering}}: میانگین هر ستون را محاسبه کرده و از تمام درایه‌های آن ستون کم می‌کنیم تا میانگین هر بعد صفر شود:
    \begin{equation}
    X_{\text{centered}} = X - \mathbf{1} \mu^T
    \end{equation}
    
    \item \textbf{ماتریس کوواریانس}: ماتریس کوواریانس نمونه $C$ را محاسبه می‌کنیم:
    \begin{equation}
    C = \frac{1}{n} X_{\text{centered}}^T X_{\text{centered}}
    \end{equation}
    
    \item \textbf{تجزیه طیفی}: مقادیر و بردارهای ویژه ماتریس کوواریانس را استخراج می‌کنیم:
    \begin{equation}
    C \cdot v_i = \lambda_i \cdot v_i
    \end{equation}
    
    \item \textbf{انتخاب مؤلفه‌های اصلی}: بردارهای ویژه $v_1, \dots, v_k$ متناظر با بزرگ‌ترین مقادیر ویژه، جهات مؤلفه‌های اصلی را مشخص می‌کنند.
\end{enumerate}
\end{algorithmbox}

\textbf{رابطه شگفت‌انگیز میان PCA و SVD}:
اگر بر روی ماتریس مرکززدایی‌شده $X_{\text{centered}}$ مستقیماً تجزیه SVD اعمال شود $(X_{\text{centered}} = U \Sigma V^T)$:
\begin{equation}
C = \frac{1}{n} (V \Sigma U^T)(U \Sigma V^T) = V \left( \frac{\Sigma^2}{n} \right) V^T
\end{equation}
این تساوی نشان می‌دهد که \textbf{بردارهای منفرد راست $V$ در SVD دقیقاً همان مؤلفه‌های اصلی در PCA هستند}، و واریانس در طول هر مؤلفه برابر با $\lambda_i = \sigma_i^2 / n$ است! در عمل، محاسبات PCA از طریق SVD انجام می‌شود تا از تشکیل ماتریس کوواریانس سنگین اجتناب گردد.

\section{تجزیه ماتریسی CUR: حفظ تنک بودن و تفسیرپذیری}
اگرچه SVD از نظر ریاضی تقریب بهینه‌ای ارائه می‌دهد، اما در پردازش کلان‌داده با دو نقص عملیاتی بزرگ همراه است:
\begin{enumerate}
    \item \textbf{تخریب ساختار تنک بودن}: ماتریس اولیه $A$ بسیار تنک است (اکثر درایه‌ها صفرند)، اما ماتریس‌های حاصل از S ($U$ و $V$) کاملاً چگال \enparen{Dense} هستند که ذخیره‌سازی آن‌ها در حافظه غیرممکن است.
    \item \textbf{فقدان تفسیرپذیری فیزیکی}: سطرهای $U$ و $V$ ترکیبات خطی عجیبی از کلمات یا کاربران هستند؛ در حالی که متخصصان داده تمایل دارند داده‌ها را بر حسب سطرها و ستون‌های واقعی توصیف کنند.
\end{enumerate}

روش \textbf{CUR}\footnote{\lr{CUR Matrix Decomposition}} برای غلبه بر این دو مشکل طراحی شد:
\begin{equation}
A \approx C \cdot U \cdot R
\end{equation}
که در آن:
\begin{itemize}
    \item $C$: شامل $c$ ستون \textbf{واقعی} انتخاب‌شده از خود ماتریس $A$ است.
    \item $R$: شامل $r$ سطر \textbf{واقعی} انتخاب‌شده از خود ماتریس $A$ است.
    \item $U$: یک ماتریس میانی کوچک با ابعاد $c \times r$ است که به عنوان پل میان $C$ و $R$ عمل می‌کند.
\end{itemize}

\subsection{الگوریتم انتخاب سطرها و ستون‌ها بر مبنای امتیاز اهرمی}
انتخاب ستون‌ها در CUR به صورت تصادفی ساده نیست، بلکه با احتمالی متناسب با مربع نرم ستون انجام می‌شود:

\begin{algorithmbox}[انتخاب ستون‌ها و سطرها در تجزیه CUR]
\begin{enumerate}
    \item به ازای هر ستون $j$ از ماتریس $A$، احتمال انتخاب را متناسب با انرژی فروبنیوس آن ستون تعریف می‌کنیم:
    \begin{equation}
    p_j = \frac{\sum_{i} A_{ij}^2}{\|A\|_F^2}
    \end{equation}
    \item تعداد $c$ ستون را با جایگذاری بر اساس توزیع احتمال $p_j$ انتخاب کرده و در ماتریس $C$ قرار می‌دهیم.
    \item برای جلوگیری از اریبی آماری، هر ستون انتخاب‌شده $j$ در ماتریس $C$ بر ضریب $\sqrt{c \cdot p_j}$ تقسیم می‌شود.
    \item به طور کاملاً متقارن، $r$ سطر از $A$ با احتمالات متناسب با انرژی سطری انتخاب شده و در ماتریس $R$ قرار می‌گیرند و بر $\sqrt{r \cdot p_i}$ تقسیم می‌شوند.
    \item \textbf{محاسبه ماتریس میانی $U$}: فرض کنید ماتریس $W$ با ابعاد $c \times r$ اشتراک سطرها و ستون‌های انتخاب‌شده در ماتریس اصلی باشد. ماتریس $U$ برابر با شبه‌وارون مور-پنروز ماتریس $W$ تعریف می‌شود:
    \begin{equation}
    U = W^+
    \end{equation}
\end{enumerate}
\end{algorithmbox}

\section{تصویرسازی‌های تصادفی و لم جانسون-لیندن‌اشتراوس \texorpdfstring{\enparen{JL Lemma}}{(JL Lemma)}}
اگر مجموعه‌ای از $n$ نقطه در یک فضای بی‌نهایت با بعد بالا $\mathbb{R}^d$ داشته باشیم، آیا می‌توان بدون از دست دادن فواصل میان نقاط، بعد فضا را کاهش داد؟

\begin{theorem}[لم جانسون-لیندن‌اشتراوس \enparen{Johnson-Lindenstrauss Lemma}]
برای هر مجموعه شامل $n$ نقطه در فضای اقلیدسی $\mathbb{R}^d$ و هر خطای مجاز $0 < \epsilon < 1$، یک نگاشت خطی تصادفی $f: \mathbb{R}^d \to \mathbb{R}^k$ وجود دارد، به گونه‌ای که با بعد هدف:
\begin{equation}
k = O\left( \frac{\log n}{\epsilon^2} \right)
\end{equation}
فواصل اقلیدسی میان \textbf{تمامی جفت‌نقاط} با تقریب $1 \pm \epsilon$ حفظ می‌شود:
\begin{equation}
(1 - \epsilon) \|x - y\|^2 \le \|f(x) - f(y)\|^2 \le (1 + \epsilon) \|x - y\|^2 \quad \forall x, y \in X
\end{equation}
\end{theorem}

\begin{examnote}
نکته اعجاب‌انگیز لم JL که پای ثابت سؤالات امتحانی است:
بعد هدف $k$ \textbf{هیچ ارتباطی با بعد اولیه $d$ ندارد}!
چه داده‌ها در فضای ۱۰۰ بعدی باشند و چه در فضای یک میلیون بعدی، اگر $n = 10^4$ نقطه داشته باشیم و خطای ۱۰٪ بخواهیم، نگاشت به فضایی با حدود چند صد بعد، تمامی فواصل را تضمین‌شده حفظ می‌کند.
ماتریس تصویرسازی تصادفی به سادگی با کشیدن درایه‌ها از یک توزیع نرمال استاندارد گاوسی $\mathcal{N}(0, 1)$ یا متغیرهای برنولی متعادل $(\pm 1)$ تولید می‌شود.
\end{examnote}

\begin{summarybox}[جمع‌بندی فصل یازدهم]
\begin{itemize}
    \item کاهش بعد راهکاری اساسی برای غلبه بر نفرین ابعاد و کشف ساختارهای پنهان در کلان‌داده است.
    \item قضیه SVD هر ماتریس را به مقادیر منفرد تجزیه کرده و طبق قضیه اکارت-یانگ بهترین تقریب با رتبه پایین را ارائه می‌دهد.
    \item الگوریتم PCA از طریق مؤلفه‌های اصلی، واریانس بیشینه داده‌های نرمال‌شده را استخراج می‌کند.
    \item تجزیه CUR با برگزیدن سطرها و ستون‌های واقعی، تنک بودن و تفسیرپذیری را حفظ می‌نماید.
    \item لم جانسون-لیندن‌اشتراوس اثبات می‌کند که تصویرسازی تصادفی به ابعادی از مرتبه $O(\frac{\log n}{\epsilon^2})$، کلیه فواصل اقلیدسی را مستقل از بعد اولیه مصون نگه می‌دارد.
\end{itemize}
\end{summarybox}
